\documentclass[12pt, a4paper]{article}
\usepackage[utf8]{inputenc}
\usepackage[T1]{fontenc}
\usepackage[french]{babel}
\usepackage{graphicx}
\usepackage{geometry}
\geometry{margin=2.5cm}
\usepackage{hyperref}
\usepackage{amsmath}
\usepackage{setspace}
\usepackage{float}
\usepackage{titlesec}

% Disable bold and italic in sections
\titleformat*{\section}{\Large}
\titleformat*{\subsection}{\large}
\titleformat*{\subsubsection}{\normalsize}

\onehalfspacing

\title{Rapport d'Analyse et de Conception : Exploitation des Données d'Accidents}
\author{DataValueXploring}
\date{\today}

\begin{document}

\maketitle
\tableofcontents
\newpage

\section{Présentation du rôle et de l'enjeu métier}

La prévention des risques routiers et l'optimisation des infrastructures de transport comptent parmi les défis majeurs des collectivités territoriales. Dans le cadre de cette étude, nous nous plaçons dans la position du Department of Transportation, l'agence de voirie américaine chargée de la gestion, de l'entretien et de l'amélioration continue du réseau routier sur un territoire donné. La sécurité des usagers est sa priorité, dans un contexte de contraintes budgétaires strictes face à un réseau extrêmement étendu.

La question centrale peut se formuler ainsi : comment réduire durablement le nombre et la gravité des accidents, en optimisant l'allocation d'un budget de travaux limité ? Aujourd'hui, la priorisation des chantiers de voirie reste largement réactive, dictée par les remontées du terrain, les plaintes de riverains ou la médiatisation de certains accidents. Ces signaux sont utiles mais partiels : ils ne suffisent pas à construire une stratégie proactive et cohérente à l'échelle d'un territoire. Une approche fondée sur les données permet de dépasser ce biais de médiatisation pour se concentrer sur le risque structurel réel.

Notre mission consiste à fournir une priorisation fondée sur des données tangibles, explicable auprès des élus et directement exploitable par les ingénieurs de la voirie. Le décideur public a besoin de réponses concrètes à deux questions : où intervenir en priorité, et quel aménagement installer une fois sur place pour réduire effectivement le risque. Il s'agit donc de transformer un volume massif de données historiques d'accidents en un plan d'action prescriptif — non de dresser un simple constat, mais d'identifier des leviers d'action infrastructurels.

% CORRECTION : « sept millions » -> chiffre exact (7,73 M)
Les données mobilisées proviennent du jeu de données US Accidents, qui couvre la période de février 2016 à mars 2023 et rassemble près de huit millions d'enregistrements (7,73 millions) répartis sur quarante-neuf États, avec de nombreuses variables décrivant le contexte météorologique, temporel et infrastructurel de chaque accident. Cette richesse d'information permet de faire évoluer la gestion de la sécurité routière d'une logique empirique vers une logique fondée sur la preuve.

Ce document détaille chaque étape du pipeline analytique développé, depuis la donnée brute jusqu'aux recommandations localisées, et justifie les choix méthodologiques retenus.

\section{Reformulation du problème en termes de données et d'apprentissage automatique}

% CORRECTION coquille : « Quand. qui appellent » -> « Quand, qui appellent »
La problématique métier, si elle est claire pour un décideur, ne se résout pas directement par un algorithme : elle doit être déconstruite en tâches spécifiques d'analyse de données et d'apprentissage automatique. Nous avons choisi de structurer la réponse autour de trois questions : Où, Quoi, Quand, qui appellent chacune des outils méthodologiques distincts.

Le premier axe, le \og Où \fg{}, consiste à identifier les lieux qui concentrent le plus grand risque d'accidents graves. Il s'agit d'une tâche d'agrégation géospatiale : non pas prédire le lieu du prochain accident de façon stochastique, mais calculer une densité d'accidents historiques normalisée par un indicateur d'exposition. Concrètement, cela revient à discrétiser l'espace en cellules homogènes, à y agréger les événements, puis à diviser cette somme par un indice de fréquentation du réseau. L'objectif est d'éviter l'écueil classique consistant à confondre zones dangereuses et zones simplement très fréquentées, où le nombre brut d'accidents est mécaniquement élevé.

Le deuxième axe, le \og Quoi \fg{}, porte sur le choix de l'aménagement une fois une zone à risque identifiée : faut-il un feu de signalisation, un panneau stop, un rond-point ? Nous formulons cette question comme un problème de classification supervisée, mais employé de manière volontairement détournée. Plutôt que de prédire l'avenir, le modèle est entraîné à prédire la gravité d'un accident historique à partir des caractéristiques de l'infrastructure et du contexte. Son rôle n'est pas de deviner, mais de servir de moteur de preuve : en interrogeant son fonctionnement interne au moyen de méthodes d'explicabilité, on identifie les éléments qui aggravent ou atténuent la probabilité de gravité. C'est cette interprétation post-hoc qui fournit la prescription, non la prédiction elle-même.

Le troisième axe, le \og Quand \fg{}, concerne les risques liés non pas à des défauts d'infrastructure permanents, mais à des conditions météorologiques ou temporelles spécifiques. Il s'agit ici d'une analyse descriptive fine des patterns de gravité selon l'heure, le jour de la semaine ou les conditions atmosphériques, afin d'orienter vers des mesures ponctuelles — panneaux à messages variables, par exemple — plutôt que vers des travaux permanents.

Cette reformulation fait ressortir un point essentiel : la performance prédictive brute (accuracy, aire sous la courbe ROC) n'est pas notre objectif. Un modèle très précis mais totalement opaque serait inutilisable pour une agence de voirie, faute d'indiquer les actions à entreprendre. Nous privilégions donc des approches où l'explicabilité et la quantification de l'effet des variables priment sur la seule performance : le Machine Learning est ici un outil de découverte de connaissances et d'évaluation de politiques publiques, plus qu'un générateur de prévisions.

\section{Description de la démarche suivie}

La démarche suit un flux de travail systématique, de la préparation des données brutes à la synthèse décisionnelle, conçu pour rester reproductible sur n'importe quel sous-ensemble territorial du jeu de données. La qualité hétérogène des données initiales a exigé une attention particulière lors des phases de nettoyage.

\subsection{Collecte, nettoyage et fiabilisation des données}

% CORRECTIONS : imputation « par localité » -> médiane ; « coordonnées aberrantes /
% hors frontières » -> exclusion des enregistrements sans coordonnées ou sans gravité
% (c'est ce que fait réellement le code, et c'est vérifié par des assertions).
La première étape a consisté à ingérer et traiter le volumineux jeu de données US Accidents. Issues de capteurs de trafic, d'applications de navigation ou de rapports de police, ces données sont par nature imparfaites et parfois inconsistantes. Un pipeline de nettoyage défensif a été développé : les valeurs manquantes des variables météorologiques (température, humidité, visibilité) ont été imputées par la médiane — un choix robuste aux valeurs extrêmes —, les colonnes trop lacunaires pour être fiables (coordonnées de fin d'accident, manquantes à 44~\%) ont été abandonnées, et les enregistrements sans coordonnées ou sans gravité renseignée ont été systématiquement exclus. L'intégrité du résultat est verrouillée par des tests automatisés exécutés à chaque passage du pipeline.

Une décision structurante a été de redéfinir l'échelle de gravité. Le jeu de données propose une échelle de un à quatre, dépendante de critères souvent subjectifs des agences de signalement locales. Pour permettre une décision binaire et pragmatique — faut-il financer une intervention urgente ou non — nous avons regroupé ces classes en une variable cible binaire : accident grave contre accident non grave. Ce choix simplifie la modélisation tout en restant aligné sur les priorités du décideur public, centrées sur la réduction de la mortalité et des blessures sévères.

Nous avons ensuite extrait et encodé les variables temporelles pertinentes (heure, jour de la semaine, mois, jours fériés) ainsi que les variables d'infrastructure au point GPS de l'accident (présence de ronds-points, de feux tricolores, de ralentisseurs, de passages à niveau). Ce travail d'ingénierie des caractéristiques est déterminant, puisque ce sont ces variables actionnables qui alimenteront le modèle.

\subsection{Analyse exploratoire temporelle et conjoncturelle (le Quand)}

L'analyse exploratoire a d'abord porté sur les dimensions temporelles et météorologiques, afin d'isoler les facteurs conjoncturels — ceux que l'aménagement urbain ne peut pas traiter directement.

\begin{figure}[H]
\centering
\includegraphics[width=0.8\textwidth]{impact_evenements.png}
\caption{Impact des événements temporels et de la météo sur la proportion de gravité.}
\end{figure}

% PRÉCISION chiffrée (vérifiée sur les 7,73 M de lignes) : 4 h -> 21,4 % de graves,
% après-midi (14 h) -> 18,2 %. Éviter les chiffres non vérifiés.
Les résultats sont par endroits contre-intuitifs. La proportion d'accidents graves n'augmente pas nécessairement lors des grands jours fériés, malgré une perception publique centrée sur les comportements à risque de ces périodes. De même, les conditions météorologiques très dégradées (brouillard épais, fortes pluies) incitent paradoxalement les conducteurs à réduire leur vitesse, ce qui se traduit par une baisse de la proportion d'accidents graves — même si le nombre absolu d'accrochages matériels peut, lui, augmenter. À l'inverse, les heures creuses affichent des taux de gravité nettement plus élevés : la part d'accidents graves culmine à 21,4~\% à quatre heures du matin, contre 18,2~\% au plus fort de l'après-midi — l'absence de trafic favorise des vitesses excessives qui transforment la moindre erreur de conduite en accident grave.

Ces constats permettent d'écarter les facteurs conjoncturels de la priorisation infrastructurelle. Les risques liés aux variations météorologiques ou aux comportements nocturnes relèvent de la régulation dynamique du trafic (signalisation d'information, limitation de vitesse temporaire, éclairage), de campagnes de prévention ou de contrôles ciblés — non de travaux de voirie. Le budget de rénovation, lui, doit se concentrer sur le risque structurel.

\subsection{Identification des zones prioritaires (le Où)}

Discrétiser l'espace géographique de façon rigoureuse était une étape indispensable. Les limites administratives (villes, comtés, États) ne conviennent pas : leurs tailles et formes hétérogènes faussent les densités calculées (effet d'unité spatiale modifiable). Nous avons retenu le système d'indexation spatiale hexagonale hiérarchique H3, qui découpe le territoire en un pavage d'hexagones réguliers. À la résolution choisie, chaque hexagone représente environ 0,7 km², soit l'échelle d'un grand carrefour, d'un rond-point majeur ou d'un tronçon de boulevard significatif.

% CORRECTION : description alignée sur le code réel — classement par charge pondérée
% gravité, puis normalisation par l'exposition OSM sur les zones de tête (pas de
% pondération par type de route).
Le nombre brut d'accidents agrégé par hexagone est cependant trompeur, à double titre. D'abord, tous les accidents ne se valent pas : les zones sont classées par leur charge, c'est-à-dire la somme des sévérités de leurs accidents, afin qu'une zone à cent accidents graves pèse davantage qu'une zone à cent accrochages bénins. Ensuite, un carrefour de centre-ville affichera toujours plus d'accidents qu'une intersection isolée en zone rurale, simplement parce qu'il voit passer un volume de véhicules bien plus important. Pour corriger ce biais d'exposition, nous avons intégré des données OpenStreetMap donnant la longueur du réseau routier contenu dans chaque hexagone de tête de classement, et calculé l'intensité d'accidents par kilomètre de voirie — un indicateur qui distingue les zones réellement dangereuses des zones simplement très pourvues en routes.

Cette démarche a permis d'établir un classement objectif et défendable des zones, et d'isoler les cinquante zones les plus critiques. Ces points noirs, bien que représentant une fraction infime du territoire, concentrent une part disproportionnée du fardeau accidentogène — un ciblage qui constitue la première brique de la rationalisation budgétaire.

\subsection{Modélisation de la gravité et identification des leviers (le Quoi)}

Pour les cinquante zones prioritaires, il restait à déterminer précisément quels équipements faisaient défaut, ou à l'inverse, quels équipements existants étaient associés à une surmortalité.

Nous avons entraîné plusieurs modèles de classification supervisée pour prédire la variable cible de gravité à partir des caractéristiques locales (feu de signalisation, passage piéton, panneau stop, ralentisseur), en incluant systématiquement des variables de contrôle — conditions météorologiques, luminosité, heure, type de route — afin d'isoler l'effet propre de l'infrastructure.

% PRÉCISION : ajout du gradient boosting (c'est le modèle effectivement retenu).
Une régression logistique et des modèles non linéaires (arbres de décision, forêts aléatoires, gradient boosting) ont été comparés méthodiquement, chacun étant en outre confronté à une base de référence naïve qui situe le niveau \og ne rien apprendre \fg{}. Une fois le meilleur modèle sélectionné et validé sur un jeu de test à découpage temporel, sur la base de métriques adaptées aux données déséquilibrées (aire sous la courbe précision-rappel), nous avons détourné sa fonction première : au lieu de prédire de nouveaux accidents, nous l'avons exploité de manière contrefactuelle via une méthode d'inférence causale, le calcul-G.

Pour chaque accident de l'échantillon, le modèle prédit d'abord la probabilité d'issue grave dans les conditions réelles. Nous modifions ensuite artificiellement une seule variable — par exemple, la présence d'un panneau stop là où il n'y en avait pas — en maintenant toutes les autres conditions identiques, puis demandons au modèle de recalculer la probabilité. L'écart moyen entre ces deux estimations, sur un large échantillon, donne l'effet ajusté de l'équipement : c'est cet effet, débarrassé de l'influence du contexte, qui fonde la recommandation d'aménagement.

\section{Justification des choix méthodologiques}

Chaque décision technique découle des contraintes des données réelles et des besoins du décideur public, et doit rester défendable devant un panel d'élus.

\subsection{Le choix de l'indexation hexagonale et de la normalisation de l'exposition}

La grille H3 se justifie par la nécessité d'unités d'analyse homogènes et équidistantes. Contrairement aux grilles rectangulaires, qui souffrent de distorsions aux pôles et aux bords, le pavage hexagonal garantit une distance identique entre le centre de chaque cellule et celui de ses voisines, ce qui facilite les calculs de voisinage et assure une représentation non biaisée du territoire.

L'intégration des données OpenStreetMap pour calculer l'exposition kilométrique est le choix qui évite le principal écueil de cette approche : sans normalisation, le modèle désignerait invariablement les grands boulevards et nœuds autoroutiers embouteillés, où les accrochages à faible vitesse sont nombreux mais rarement graves — au détriment de routes secondaires moins fréquentées mais proportionnellement plus meurtrières du fait de leur configuration.

Le kilométrage OpenStreetMap reste un proxy imparfait du trafic réel (nous ne disposons pas de comptages exhaustifs à l'échelle nationale), mais c'est le plus robuste et le plus disponible pour garantir la reproductibilité de la méthode sur n'importe quel État.

% CORRECTION : légende alignée sur le contenu réel de la figure (distribution des
% scores + relation fréquence/gravité), qui montre la concentration extrême du risque.
\begin{figure}[H]
\centering
\includegraphics[width=0.8\textwidth]{distribution_scores.png}
\caption{Concentration extrême du risque : distribution des scores des cellules H3 (une infime fraction des zones au-dessus du 99\textsuperscript{e} centile) et relation entre fréquence et gravité moyenne.}
\end{figure}

La normalisation par l'exposition fait ensuite chuter les scores des zones très étendues ou très denses et fait émerger les véritables points noirs : celles où le rapport entre nombre d'accidents et longueur de voirie disponible est anormalement élevé.

\subsection{Le modèle d'apprentissage comme moteur d'explication}

Le Machine Learning est le plus souvent déployé comme une science de prédiction pure. Dans ce contexte, la prédiction pure est une impasse : un modèle boîte noire annonçant avec certitude qu'un accident grave surviendra demain à tel croisement n'aide en rien à le prévenir s'il ne dit pas pourquoi ce croisement est dangereux.

% PRÉCISION : gradient boosting ajouté ici aussi.
C'est pour répondre à ce besoin de prescriptivité que nous avons retenu une démarche explicative post-modélisation. Comparer une régression logistique simple à des modèles ensemblistes (Random Forest, gradient boosting) répond à un besoin de vérification croisée.

La régression logistique s'interprète directement via ses coefficients, traduits en rapports de cotes (odds ratios), qui indiquent de combien la probabilité d'accident grave est multipliée en présence d'un élément d'infrastructure donné. Elle impose toutefois une hypothèse de linéarité qui peut s'avérer simpliste en environnement urbain complexe. Nous entraînons donc en parallèle des modèles ensemblistes, mieux à même de capturer les non-linéarités, et utilisons des techniques d'explicabilité indépendantes du modèle — les valeurs de Shapley — pour vérifier si les relations identifiées par la régression logistique se confirment dans l'espace non linéaire des forêts aléatoires. La convergence de ces deux lectures constitue notre garantie de validité.

\subsection{La subtilité des valeurs d'explicabilité et les biais de confusion}

L'analyse des valeurs de Shapley a d'abord suggéré, en effet global, que la présence de feux de signalisation et de passages piétons augmentait la probabilité de gravité. Une lecture non critique de ce résultat conduirait à recommander leur suppression — une conclusion absurde et dangereuse.

Cette apparente anomalie s'explique par un biais de sélection classique : feux et passages piétons ne sont pas installés au hasard, mais dans les zones urbaines les plus denses et les intersections les plus complexes, précisément là où le risque intrinsèque — lié à la cohabitation de véhicules, cyclistes et piétons — est le plus élevé. L'équipement signale le danger et tente de le réguler ; il ne le crée pas.

% CORRECTION : description alignée sur le code réel — l'effet conditionnel est mesuré
% sur les accidents survenus en présence de l'aléa local (pas une stratification par
% « zones de complexité comparable »), le contexte étant contrôlé par le modèle.
C'est pourquoi nous n'utilisons jamais les effets globaux bruts pour orienter une décision d'infrastructure. Pour neutraliser ce biais, nous calculons des effets ajustés conditionnels : l'impact de chaque aménagement (via le calcul-G) est mesuré uniquement sur les accidents survenus en présence de l'aléa local dominant de la zone considérée, le contexte — météo, heure, luminosité, État — étant maintenu constant par le modèle. Ce conditionnement neutralise le facteur de confusion lié à la densité urbaine et révèle l'effet protecteur réel des équipements.

\subsection{La validation empirique de la stabilité temporelle du risque}

Les travaux d'infrastructure se planifient et s'exécutent sur plusieurs années ; encore faut-il s'assurer que le risque identifié est structurel, et non lié à un chantier temporaire ou à une anomalie statistique ponctuelle — sans quoi l'investissement serait un gaspillage de fonds publics.

\begin{figure}[H]
\centering
\includegraphics[width=0.8\textwidth]{validation_stabilite.png}
\caption{Validation empirique de la stabilité temporelle des cinquante zones prioritaires à travers deux périodes distinctes.}
\end{figure}

Pour le vérifier, nous avons scindé les données en deux périodes chronologiques distinctes (2016--2019 et 2020--2023) et réexécuté indépendamment le pipeline de scoring géographique sur chacune. Nous avons ensuite mesuré la corrélation de rang de Spearman entre les cinquante zones les plus dangereuses des deux périodes.

Les résultats montrent une corrélation élevée et une stabilité marquée : les zones prioritaires ne résultent pas d'un bruit statistique, mais d'un dysfonctionnement structurel persistant du réseau — un constat qui donne au décideur politique la garantie nécessaire pour justifier un investissement pérenne.

% ============================================================================
% NOUVELLE SECTION — pistes explorées, impasses et difficultés
% ============================================================================
\section{Pistes explorées et difficultés rencontrées}

La solution présentée dans les sections précédentes est le produit d'une démarche itérative : plusieurs pistes ont été ouvertes, certaines abandonnées, et deux erreurs méthodologiques significatives ont été commises puis corrigées. Nous les documentons ici, car elles ont directement façonné les garde-fous du livrable final.

\subsection{Pistes ouvertes puis écartées}

Trois pistes envisagées lors de l'analyse préliminaire n'ont pas été retenues. La première était un modèle de \emph{durée de blocage} du trafic, exploitant l'écart entre l'heure de début et de fin de chaque accident ; elle a été abandonnée lorsque l'exploration a montré que 44~\% des coordonnées de fin étaient manquantes, jetant un doute sérieux sur la fiabilité des horodatages de fin eux-mêmes. La deuxième concernait le périmètre géographique : nous avons hésité entre restreindre l'étude au comté de Los Angeles (lisibilité maximale), à un État médian comme la Caroline du Sud (volume modeste mais équilibré), ou à la Californie entière. Le débat a été tranché par un argument de généricité : entraîner sur l'ensemble des États-Unis et paramétrer la démonstration par une simple constante de territoire, la Californie ne servant que de terrain d'illustration. La troisième piste était une chaîne de modélisation entièrement portée par Spark~MLlib ; elle a été écartée au profit de scikit-learn sur un échantillon plafonné et documenté, seul moyen de mener dans le temps imparti une comparaison riche de plusieurs modèles avec les métriques et les outils d'explicabilité souhaités. MLlib reste mentionné comme voie de passage à l'échelle.

\subsection{Une impasse instructive : le modèle tautologique}

La première tentative de réponse au \og Quoi \fg{} constitue l'erreur la plus instructive du projet. Nous avions entraîné un modèle de régression (forêt aléatoire) prédisant la \emph{charge} d'une zone — la somme des sévérités de ses accidents — à partir des comptages d'accidents impliquant chaque équipement dans cette même zone. Le coefficient de détermination était excellent, et pour cause : la cible était une fonction quasi arithmétique des variables d'entrée, puisque les comptages par équipement \emph{composent} la charge. Le modèle n'apprenait pas un phénomène routier ; il reconstituait une identité comptable, et ses importances de variables ne décrivaient que la composition du territoire.

Cette impasse a produit deux garde-fous permanents. D'une part, le changement d'unité d'analyse : le modèle final travaille au niveau de l'accident individuel, dont la gravité n'est pas une fonction de ses caractéristiques d'infrastructure. D'autre part, des tests automatisés anti-fuite, exécutés à chaque exécution du notebook, qui vérifient qu'aucune variable d'entrée ne contient la cible sous quelque forme que ce soit et que le découpage temporel entre entraînement et test est étanche.

\subsection{Le piège des effets bruts, en deux temps}

La seconde correction majeure concerne la chaîne d'explication elle-même, et s'est jouée en deux temps. Le premier est décrit en section~4.3 : les valeurs de Shapley globales suggéraient que feux et passages piétons aggravaient la gravité, artefact du biais de sélection des équipements. Le second est plus subtil : notre première version de la prescription par zone reposait sur des écarts de sévérité moyens calculés directement dans les données — alors même que l'argument central du recours au modèle était sa capacité à contrôler les facteurs de confusion. La prescription contredisait donc silencieusement sa propre justification. Nous avons résolu cette incohérence en faisant piloter la recommandation par l'effet ajusté conditionnel issu du calcul-G, l'écart brut observé n'étant conservé que comme contre-vérification. C'est également ce qui a motivé l'ajout des rapports de cotes de la régression logistique, avec leurs intervalles de confiance : deux lectures indépendantes qui concordent sur dix équipements sur douze constituent une garantie qu'aucune des deux ne fournit isolément.

\subsection{Difficultés d'ingénierie}

Deux difficultés techniques méritent mention. La première est la volumétrie : avec près de huit millions de lignes, chaque erreur de conception de la pipeline Spark se paie en dizaines de minutes de recalcul ; la discipline de mise en cache des jeux intermédiaires et de plafonnement explicite des échantillons a été acquise par la pratique. La seconde, plus sournoise, est la dépendance au service Overpass d'OpenStreetMap pour l'exposition routière : requêtes bloquantes sans délai de garde, limitations de débit agressives et imprévisibles, et un piège de configuration — la clé de cache du client dépend du délai d'expiration configuré, si bien qu'un cache constitué avec un réglage différent devient invisible. La version finale encapsule cette étape dans un budget de temps global avec dégradation contrôlée : en l'absence de réponse du service, le classement par volume est conservé et la limite documentée, de sorte que l'exécution complète du notebook aboutit dans tous les cas. Le cache local est en outre livré avec le dépôt, ce qui permet de rejouer l'exposition hors-ligne sur quarante-deux des cinquante zones.

Enfin, la suite de tests a démontré son utilité dès sa première exécution en détectant une régression de contenu : une édition antérieure du notebook avait amputé, sans que personne ne s'en aperçoive, la section d'analyse critique. L'anecdote illustre concrètement pourquoi un livrable, même sous forme de notebook, gagne à être protégé par un protocole de validation automatisé.

\section{Valeur créée et limites de l'approche}

Ce travail ne prend son sens que traduit en valeur actionnable pour l'utilisateur métier. Il est tout aussi nécessaire d'en exposer les limites avec transparence : ignorer les angles morts d'un modèle est le plus sûr moyen de compromettre son adoption.

\subsection{Valeur pour la décision stratégique et opérationnelle}

La première valeur de cette solution est d'objectiver un débat souvent politique et passionné autour de l'allocation des budgets de sécurité routière. Le département de la voirie dispose désormais d'un outil quantitatif, auditable et reproductible, qui produit une liste hiérarchisée de zones à traiter en priorité.

L'outil va cependant au-delà d'un simple classement : pour chaque zone prioritaire, il croise l'aléa local avec l'effet ajusté conditionnel des mesures protectrices actuellement absentes. Le résultat n'est pas une statistique abstraite, mais une recommandation d'aménagement lisible par un ingénieur — par exemple, l'installation d'une signalisation avancée sur un carrefour présentant un historique de collisions latérales, justifiée par un impact conditionnel négatif significatif sur la probabilité de gravité, toutes choses égales par ailleurs sur la météo ou l'éclairage.

Cette recommandation s'accompagne d'une estimation quantitative ex ante de l'impact attendu : en multipliant l'effet ajusté par le volume historique d'accidents pertinents dans la zone, le pipeline fournit une estimation du nombre d'accidents graves potentiellement évitables. Ce chiffre, même approximatif, permet de construire de véritables analyses coût-bénéfice avant l'arbitrage budgétaire — une valeur à la fois technique (quoi faire) et politique (pour quel bénéfice).

\subsection{Limites liées aux données et à la modélisation}

Plusieurs limites méritent d'être assumées clairement.

La première, la plus fondamentale, est l'absence de causalité formellement démontrée. Nos recommandations reposent sur des associations statistiques ajustées, non sur des expériences contrôlées randomisées. Le jeu de données ne renseigne pas la date d'installation des équipements : nous ne pouvons donc pas observer un même carrefour avant et après l'ajout d'un feu ou d'un rond-point, et devons nous contenter de comparaisons transversales entre carrefours équipés et non équipés. Malgré les variables de contrôle incluses, des facteurs confondants non observés (qualité du revêtement, angle de visibilité, végétation, signalisation horizontale effacée) peuvent subsister et biaiser partiellement les estimations.

La deuxième limite concerne les facteurs humains, que l'approche infrastructurelle laisse volontairement de côté. Un accident résulte presque toujours de la conjonction d'une infrastructure défaillante et d'une erreur humaine ; or le modèle ignore l'âge du conducteur, sa fatigue, une éventuelle alcoolémie ou sa vitesse réelle d'impact. Son pouvoir prédictif global reste donc structurellement plafonné. Sur certaines zones très accidentogènes, le modèle n'identifie d'ailleurs aucun levier d'aménagement pertinent — non pas un échec, mais un résultat en soi : il signale que la racine du problème relève probablement du comportement des usagers plutôt que de l'ingénierie civile, et appelle des réponses de sensibilisation ou de contrôle plutôt que des travaux.

Enfin, un biais structurel affecte la collecte elle-même : les fournisseurs de données de trafic et les applications de navigation sur-représentent les grands axes et les voies urbaines denses, et sous-déclarent probablement les accrochages mineurs sur les routes rurales isolées ou les zones résidentielles peu connectées. La carte des risques produite reste donc tributaire de la qualité de son alimentation.

\subsection{Perspectives d'évaluation empirique sur le terrain}

Pour compenser l'absence de preuve causale directe dans les données historiques, la légitimité de l'outil doit se construire par une évaluation empirique sur le terrain.

Nous recommandons de ne pas généraliser les aménagements à l'échelle de l'État de façon précipitée, mais de sélectionner un petit groupe de zones pilotes — par exemple dix intersections classées en priorité la plus élevée — où les aménagements prescrits seront effectivement réalisés. En parallèle, un groupe de contrôle de dix zones aux scores de risque et configurations comparables ne recevra aucun aménagement dans un premier temps. Un suivi sur un à deux ans, comparant l'évolution du nombre d'accidents graves entre groupe traité et groupe témoin, permettrait d'appliquer une méthode de différence de différences (Difference-in-Differences) pour quantifier l'impact causal réel des travaux.

C'est à l'issue de cette expérimentation que l'estimation ex ante du nombre d'accidents évités, calculée par le modèle, pourra être confrontée à la réalité du terrain et validera pleinement la chaîne de valeur de cette démarche.

\section{Conclusion}

Ce projet montre comment des données massives, imparfaites et hétérogènes peuvent être transformées en un levier d'action publique et de pilotage stratégique. En refusant la prédiction en boîte noire et en exigeant du modèle qu'il justifie ses liens avec l'infrastructure de façon interprétable, nous avons construit un processus décisionnel transparent et défendable — y compris en documentant les impasses rencontrées en chemin, qui ont façonné les garde-fous méthodologiques de la version finale.

L'agence de voirie dispose désormais d'un protocole complet et reproductible, de l'identification géographique d'un danger chronique jusqu'à l'estimation quantitative a priori de l'impact social d'une rénovation ciblée. Ce système n'a pas vocation à remplacer le jugement qualitatif, la connaissance locale ni l'expertise des ingénieurs de terrain ; il en rationalise le processus en amont, oriente les budgets là où l'impact potentiel sur la sécurité publique est le plus grand, et fournit une argumentation fondée sur des preuves pour la gestion de la sécurité routière.

\end{document}
